\documentclass[10pt,a4paper]{amsart}
\usepackage[margin=19mm]{geometry}
\usepackage{amsmath,amssymb,mathtools}
\usepackage[scr=rsfs,cal=euler]{mathalfa}
\usepackage{xfrac}
\usepackage[unicode,hidelinks,bookmarks=false]{hyperref}
\newcommand{\ie}{i.e.\ }
\newcommand{\cf}{cf.\ }
\newcommand{\resp}{respectively\ }
\newcommand{\Subsection}{\subsection}
\providecommand{\nobreakdash}{\nobreak\hbox{-}}
\setlength{\emergencystretch}{2em}
\multlinegap0pt
\usepackage{bm}
\usepackage{bbm}
\usepackage{dsfont}
\usepackage{xspace}
\usepackage{cancel}
\usepackage{mathtools}
\usepackage[shortlabels]{enumitem}
\usepackage{amsthm}
\makeatletter
\@ifundefined{claim}{\newtheorem{claim}{Claim}}{}
\@ifundefined{theorem}{\newtheorem{theorem}{Theorem}}{}
\@ifundefined{lemma}{\newtheorem{lemma}[theorem]{Lemma}}{}
\@ifundefined{proposition}{\newtheorem{proposition}[theorem]{Proposition}}{}
\@ifundefined{remark}{\newtheorem{remark}[theorem]{Remark}}{}
\makeatother
\title[On maximal ladders]{On maximal ladders}
\author{Lorenzo Notaro}
\date{Source: arXiv:2604.06031v1 (CC BY 4.0)}
\begin{document}
\maketitle

\noindent\emph{Source and licence.} On maximal ladders. Version arXiv:2604.06031v1 (CC BY 4.0), distributed under the Creative Commons Attribution 4.0 International licence, as recorded in the arXiv licence metadata for this exact version and shown on the abstract page https://arxiv.org/abs/2604.06031v1. Full rights notice: licenses/arxiv-2604.06031-CC-BY-4.0.txt.\par\smallskip

\noindent\textit{Editorial scope.} Counted unit: the Proposition whose source label is \texttt{prop:wehrung} in the article (source0.tex lines 378--380, source label \texttt{prop:wehrung}), delivered together with the article's own complete original proof of it (source0.tex lines 381--531, one \texttt{proof} environment of 151 lines). No mathematics is written, repaired or re-derived: statement and proof are the article's own text line for line. Required context retained uncounted: lemma:meetproj (lines 211--213); lemma:concproj (lines 222--224); thm:characterization (lines 302--304); rmk:cofinalmeet (lines 328--330). Every definition, display and label of those lines is retained unchanged. Omitted from this package and not redistributed: 1--210, 214--221, 225--301, 305--327, 331--377, 532--1617. Nothing inside a retained excerpt is omitted or altered: every retained body is delivered exactly as the article's lines are written, so no omission receipt of any kind accompanies this package. Citations: the retained excerpts carry no citation command at all, so no bibliography entry of the article is transcribed and no reference is redistributed. Reference adaptation: none was needed and none was made; every reference inside the delivered unit resolves inside it, so no unresolved reference or citation is produced. Printed slips kept literal: no printed word, symbol, hypothesis or number of the counted statement or of its proof was altered. Layout and class: the article's own class is \texttt{amsart} with packages including amsmath, amsthm, amsfonts, amssymb, enumitem, todonotes, float, mathtools, csquotes, changes, comment, tikz-cd, biblatex, hyperref; the wrapper is the lane's pinned \texttt{amsart} editorial wrapper at 10pt on A4, which restates the theorem-like environments the retained text needs and adds the article's own macro definitions that the retained text needs (none), with extra wrapper packages (bm, bbm, dsfont, xspace, cancel, mathtools, enumitem). The delivered numbering is the unit's own. Assets: no retained span contains a figure environment, a graphics-inclusion command, a PDF-page inclusion, a table or a TikZ picture; no image file is copied into this package, the package declares no asset dependency and no image file is redistributed with it. Licence: version arXiv:2604.06031v1 of this article is distributed under the Creative Commons Attribution 4.0 International licence, as recorded in the arXiv licence metadata for this exact version and shown on the abstract page https://arxiv.org/abs/2604.06031v1. A full rights notice is delivered at licenses/arxiv-2604.06031-CC-BY-4.0.txt. Characters: every retained excerpt is pure ASCII, so the delivered engine, XeLaTeX with its default Latin Modern text font, reports no missing character.
\begin{lemma}\label{lemma:meetproj}
Given a lower finite lattice $P$, some elements $x,y \in P$ and an ideal $I \subseteq P$, $\pi_I(x \wedge y) = \pi_I(x) \wedge \pi_I(y)$.
\end{lemma}

\begin{lemma}\label{lemma:concproj}
If $P$ is a lower finite lattice and $I,J \subseteq P$ are ideals with $I \subseteq J$, then $\pi_I \circ \pi_J = \pi_I$.
\end{lemma}

\begin{theorem}\label{thm:characterization}
For every $n > 1$, an $n$-ladder is maximal if and only if it has no cofinal meet-subsemilattice which is also a $(n-1)$-ladder in the induced order.
\end{theorem}

\begin{remark}\label{rmk:cofinalmeet}
As noted in the proof of Theorem~\ref{thm:characterization}, a cofinal meet-subsemilattice of a lower finite lattice is itself a lattice in the induced ordering. In particular, an $n$-ladder $(P, \le) $ is maximal if and only if it has no cofinal meet-subsemilattice $C \subseteq P$ whose elements have at most $n-1$ lower covers in $(C, \le)$.
\end{remark}

\begin{proposition}\label{prop:wehrung}
Given some $n > 0$ and a lower finite lattice $L$ of cardinality $\aleph_n$, $\text{\upshape Add}(\omega, \omega_n)$ forces the existence of a cofinal meet-subsemilattice of $L$ which is an $(n+1)$-ladder.
\end{proposition}

\begin{proof}
Fix a lower finite lattice $(L, \le)$ of cardinality $\aleph_n$. We first partition $L$ into countable pieces indexed by certain terminal nodes of a tree. Then, using this partition, we show that $\text{Add}(\omega, \omega_n)$ forces the existence of a cofinal meet-subsemilattice of $L$ which is an $(n+1)$-ladder in the induced ordering. 

Let
\begin{multline*}
T \coloneqq \big\{\emptyset\} \cup \{\langle\alpha_0, \alpha_1, \dots, \alpha_{m}\rangle : m < n \text{ and } \alpha_i < \omega_{n-i} \text{ for every } i \le m \text{ and } \\ \alpha_i \text{ is not a limit ordinal for every } i < m\big\}.
\end{multline*}
Given $\vec{\alpha}, \vec{\beta} \in T$, we say that $\vec{\beta}$ \emph{extends} $\vec{\alpha}$, in symbols $\vec{\alpha} \subseteq \vec{\beta}$, when $\vec{\beta} \upharpoonright |\vec{\alpha}| = \vec{\alpha}$. Moreover, we write $\vec{\alpha} <_{lex}^* \vec{\beta}$ to mean that $\vec{\alpha} <_{lex} \vec{\beta}$ and $\vec{\alpha} \nsubseteq \vec{\beta}$.

We call $\vec{\alpha} \in T$ \emph{nonlimit} if its last element is not a limit ordinal. Moreover, we call $\vec{\alpha} \in T$ \emph{maximal} if either it has length $n$ or its last element is a limit ordinal---i.e., if it has no proper extension in $T$. 


\begin{claim}\label{prop:wehrung:claim:1}
There is a family $\langle I_{\vec{\alpha}} : \vec{\alpha} \in T\rangle$ of ideals  of $L$ such that:
\begin{enumerate}[label={\upshape (\arabic*)}]
\item $I_{\emptyset} = L$,
\item $I_{\vec{\alpha}} = \bigcup_{\beta < \omega_{n-|\vec{\alpha}|}} I_{\vec{\alpha}^\smallfrown \beta}$ if $\vec{\alpha}$ is not maximal,
\item $I_{\vec{\alpha}^\smallfrown \beta} \subseteq I_{\vec{\alpha}^\smallfrown \gamma}$ for every $\beta \le \gamma < \omega_{n-|\vec{\alpha}|}$ if $\vec{\alpha}$ is not maximal,
\item $I_{\vec{\alpha}^\smallfrown \lambda} = \bigcup_{\beta < \lambda} I_{\vec{\alpha}^\smallfrown \beta}$ if $\lambda$ is a limit ordinal and $\vec{\alpha}$ is not maximal,
\item $I_{\vec{\alpha}} \nsubseteq \bigcup_{\vec{\beta} <_{lex}^* \vec{\alpha}} I_{\vec{\beta}}$ if $\vec{\alpha}$ is nonlimit,
\item $I_{\vec{\alpha}} \supseteq I_{\vec{\gamma}}$ for every nonlimit $\vec{\gamma} \in T$ of length $n$ such that $\vec{\gamma} <_{lex} \vec{\alpha}$ and $(I_{\vec{\alpha}} \cap I_{\vec{\gamma}}) \setminus \bigcup_{\vec{\beta} <_{lex}^* \vec{\gamma}} I_{\vec{\beta}} \neq \emptyset$,
\item $|I_{\vec{\alpha}}| = \aleph_{n-|\vec{\alpha}|}$.
\end{enumerate}
\end{claim}
\begin{proof}
We define $I_{\vec{\alpha}}$ by induction on the lexicographic ordering of the $\vec{\alpha}$s in $T$. As soon as we define $I_{\vec{\alpha}}$ for some $\vec{\alpha} \in T$ we fix an enumeration $\langle x_{\vec{\alpha}}(\delta) : \delta < \omega_{n-|\vec{\alpha}|} \rangle$ of $I_{\vec{\alpha}}$. 

First let $I_{\emptyset} = L$. Fix some nonempty $\vec{\alpha} \in T$ and let $k < n$ be such that $k+1 = |\vec{\alpha}|$. Suppose that we defined $I_{\vec{\beta}}$ for all $\vec{\beta}$ in $T$ with $\vec{\beta} <_{lex} \vec{\alpha}$ such that (3)-(7) hold below $\vec{\alpha}$---i.e., the statements (3)-(7) hold when the elements of $T$ mentioned are $<_{lex} \vec{\alpha}$. Moreover, we also assume that the following weakening of (2) holds:
\begin{enumerate}
\item[(2')] $I_{\vec{\beta}} \supseteq I_{\vec{\gamma}}$ for all $\vec{\beta},\vec{\gamma} <_{lex} \vec{\alpha}$ such that $\vec{\gamma}$ extends $\vec{\beta}$.
\end{enumerate} 
We now define $I_{\vec{\alpha}}$. If $\alpha_k$ is a limit ordinal, let $I_{\vec{\alpha}} = \bigcup_{\beta < \alpha_k} I_{(\vec{\alpha} \upharpoonright k)^\smallfrown \beta}$. Now suppose that $\alpha_k$ is not a limit ordinal. First let us fix a subset $X \subseteq I_{\vec{\alpha} \upharpoonright k}$ depending on the value of $\alpha_k$:\vspace{0.5em}

\noindent \underline{$\alpha_k = 0$}:  for every $\vec{\beta}$, $\vec{\beta} <^*_{lex} \vec{\alpha}$ if and only if $\vec{\beta} <^*_{lex} \vec{\alpha} \upharpoonright k$. By (5) and (7), $I_{\vec{\alpha} \upharpoonright k} \nsubseteq \bigcup_{\vec{\beta} <^*_{lex} \vec{\alpha}}I_{\vec{\beta}}$, and $|I_{\vec{\alpha} \upharpoonright k}| = \aleph_{n-k}$. Thus, we can pick $X\subseteq I_{\vec{\alpha} \upharpoonright k}$ such that $|X| = \aleph_{n-k-1}$ and  $X \nsubseteq \bigcup_{\vec{\beta} <^*_{lex} \vec{\alpha}} I_{\vec{\beta}}$. \vspace{0.5em}

\noindent \underline{$\alpha_k > 0$}: it follows that 
\[
\{\vec{\beta} \in T : \vec{\beta} <^*_{lex} \vec{\alpha}\} = \{(\vec{\alpha} \upharpoonright k)^\smallfrown \gamma : \gamma < \alpha_k\} \cup \{\vec{\beta} \in T : \vec{\beta} <^*_{lex} \vec{\alpha}\upharpoonright k\}.
\]
 By (3), $I_{(\vec{\alpha} \upharpoonright k)^\smallfrown \gamma} \subseteq I_{(\vec{\alpha} \upharpoonright k)^\smallfrown (\alpha_k-1)}$ for every $\gamma < \alpha_k$. By (2') and (7), $I_{\vec{\alpha} \upharpoonright k} \supseteq I_{(\vec{\alpha} \upharpoonright k)^\smallfrown (\alpha_k - 1)}$ and $|I_{\vec{\alpha} \upharpoonright k}| > |I_{(\vec{\alpha} \upharpoonright k)^\smallfrown (\alpha_k - 1)}|$. Moreover, by (5), $I_{\vec{\alpha} \upharpoonright k} \nsubseteq \bigcup_{\vec{\beta} <^*_{lex} \vec{\alpha} \upharpoonright k} I_{\vec{\beta}}$. The set $I_{\vec{\alpha} \upharpoonright k} \setminus \bigcup_{\vec{\beta} <^*_{lex} \vec{\alpha} \upharpoonright k} I_{\vec{\beta}}$, being nonempty and upward closed, must have the same cardinality as $I_{\vec{\alpha} \upharpoonright k}$, that is $\aleph_{n-k}$. Indeed, in general, an infinite lower finite poset has the same cardinality as all its cofinal subsets. Thus, we can fix some $z$ belonging to the set
\[
I_{\vec{\alpha} \upharpoonright k} \setminus \Big ( I_{(\vec{\alpha} \upharpoonright k)^\smallfrown (\alpha_k - 1)} \cup \bigcup_{\vec{\beta} <^*_{lex} \vec{\alpha} \upharpoonright k} I_{\vec{\beta}} \Big).
\] 
Now, let $\delta < \omega_{n-k}$ be the least ordinal such that $
x_{\vec{\alpha} \upharpoonright k}(\delta) \not\in I_{(\vec{\alpha} \upharpoonright k)^\smallfrown (\alpha_k - 1)}$.
Set $X = I_{(\vec{\alpha} \upharpoonright k)^\smallfrown (\alpha_k - 1)} \cup \{z, x_{\vec{\alpha} \upharpoonright k}(\delta)\}$.\vspace{0.5em}


 Now that we have fixed $X$, let us define a $\subseteq$-increasing sequence $\langle F_m : m\in\omega\rangle$ of ideals of $L$ as follows: first let $F_0 = \mathrm{id}(X)$,  that is, $F_0$ is the ideal generated by $X$; now, if $F_m$ is defined, let $R_m$ be the set of all those $\vec{\gamma} \in T$ of length $n$ such that $\vec{\gamma} <_{lex} \vec{\alpha}$ and $(F_m \cap I_{\vec{\gamma}}) \setminus \bigcup_{\vec{\beta} <^*_{lex} \vec{\gamma}} I_{\vec{\beta}} \neq \emptyset$; let $F_{m+1} = \mathrm{id}(F_m \cup \bigcup_{\vec{\gamma} \in R_m} I_{\vec{\gamma}})$. Finally, let $I_{\vec{\alpha}} = \bigcup_{m\in\omega} F_m$. 
 
This ends the definition of $I_{\vec{\alpha}}$. Let us check that (2') and (3)-(7) still hold when all the members of $T$ mentioned are $\le_{lex} \vec{\alpha}$. Properties (2') and (3)-(6) follow directly from the construction of $I_{\vec{\alpha}}$. 

We are left to prove that (7) holds. It suffices to show $|I_{\vec{\alpha}}| = \aleph_{n-|\vec{\alpha}|} = \aleph_{n-k-1}$. If $\alpha_k$ is a limit ordinal, then it follows directly that $|I_{\vec{\alpha}}| = |\alpha_k| \cdot \sup_{\beta < \alpha_k} |I_{(\vec{\alpha} \upharpoonright k)^\smallfrown \beta}| = |\alpha_k| \cdot \aleph_{n-k-1}$, but since $\alpha_k < \omega_{n-k}$, we conclude $|I_{\vec{\alpha}}| = \aleph_{n-k-1}$. Now suppose that $\alpha_k$ is not a limit ordinal. It suffices to prove that $|F_m| = \aleph_{n-k-1}$ for each $m \in \omega$. Since $|X| = \aleph_{n-k-1}$ by choice and $L$ is lower finite, it follows that $|F_0| = \aleph_{n-k-1}$. Now, if $|F_m| = \aleph_{n-k-1}$, in order to prove that $|F_{m+1}| = |F_m|$ it suffices to show that $|R_m| \le |F_m|$. Indeed, it is clear that for each $x \in L$ there exists at most one $\vec{\gamma} \in T$ of length $n$ such that $x \in I_{\vec{\gamma}} \setminus \bigcup_{\vec{\beta} <^*_{lex} \vec{\gamma}} I_{\vec{\beta}}$. Then, it follows directly that $|R_m| \le |F_m|$. 

This ends the inductive definition of the $I_{\vec{\alpha}}$s. Property (1) holds by construction and we have shown that (3)-(7) also hold. Regarding (2), note that it holds by the way we chose $\delta$ in the case where $\alpha_k$ is a successor ordinal. 
\end{proof}

Fix a family satisfying the statement of Claim~\ref{prop:wehrung:claim:1}.
Given some $\vec{\alpha} \in T$, for notational clarity let us denote $\pi_{I_{\vec{\alpha}}}$ by $\pi_{\vec{\alpha}}$. Let $T^s(n)$ be the set of all the nonlimit elements of $T$ which have length $n$. For each $\vec{\alpha} \in T^s(n)$, let
\[
J_{\vec{\alpha}} \coloneqq I_{\vec{\alpha}} \setminus \bigcup_{\vec{\beta} <^*_{lex} \vec{\alpha}}I_{\vec{\beta}}.
\]
By (5) and (7), the $J_{\vec{\alpha}}$s are nonempty, countable, and pairwise disjoint.

By (1) and (2), for every $x \in L$ there exists $\vec{\alpha} \in T$ of length $n$ such that $x \in I_{\vec{\alpha}}$. Let $\vec{\alpha}(x)$ be the $\le_{lex}$-least $\vec{\alpha} \in T$ of length $n$ such that $x \in I_{\vec{\alpha}}$. By (4), $\vec{\alpha}(x) \in T^s(n)$. Also,  clearly $x \in J_{\vec{\alpha}(x)}$. In particular, we have
\begin{equation}\label{eq:wehrung}
L = \bigcup \big\{J_{\vec{\alpha}} : \vec{\alpha} \in T^s(n)\big\},
\end{equation}
that is to say, the $J_{\vec{\alpha}}$s form a partition of $L$.

Notice that for every $x,y \in L$, if $x \le y$, then $\vec{\alpha}(x) \le_{lex}\vec{\alpha}(y)$. Finally, observe that, by (6), for every $x,y \in L$, if $x \le y$, then $I_{\vec{\alpha}(x)} \subseteq I_{\vec{\alpha}(y)}$.

Consider the forcing notion $\mathbb{P}$ defined as
\[
\mathbb{P} \coloneqq \big\{p:D \rightarrow L : D \subseteq T^s(n) \times \omega \text{ is finite and } \forall (\vec{\alpha}, m) \in D \ (p(\vec{\alpha}, m) \in J_{\vec{\alpha}})\big\},
\]
and ordered by reverse extension: $q \le p$ if $\mathrm{dom}(q) \supseteq \mathrm{dom}(p)$ and $q \upharpoonright \mathrm{dom}(p) = p$. Since $T^s(n)$ has cardinality $\aleph_n$ and each $J_{\vec{\alpha}}$ is countable, it is easy to see that $\mathbb{P} \cong \text{Add}(\omega, \omega_n)$, i.e., $\mathbb{P}$ is isomorphic to Cohen's forcing $\text{Add}(\omega, \omega_n)$.

For each $p \in \mathbb{P}$, we now define a subset $C_p \subseteq L$. Fix $p \in \mathbb{P}$. By \eqref{eq:wehrung}, it suffices to define $C_p \cap J_{\vec{\alpha}}$ for each  $\vec{\alpha} \in T^s(n)$. We do so by induction on the lexicographic ordering of $T^s(n)$. 

Fix some $\vec{\alpha} \in T^s(n)$ and suppose that $C_p \cap J_{\vec{\beta}}$ is defined for every $\vec{\beta} \in T^s(n)$ with $\vec{\beta} <_{lex} \vec{\alpha}$, towards defining $C_p \cap J_{\vec{\alpha}}$. For every $x \in J_{\vec{\alpha}}$, we let $x \in C_p$ if and only if there exists some $m \in \omega$ such that:
\begin{enumerate}[label=(\alph*)]
\item $p (\vec{\alpha}, m) = x$, and
\item for all $k < m$,  $(\vec{\alpha}, k) \in \mathrm{dom}(p)$ and $p(\vec{\alpha}, k) \le x$, and
\item for all $\vec{\gamma} \in T$ such that $\vec{\gamma} <_{lex}^* \vec{\alpha}$ ,  $\pi_{\vec{\gamma}}(x) \in C_p$.
\end{enumerate}
Note that if $p,q \in \mathbb{P}$ and $p \le q$, then $C_q \subseteq C_p$. Now, given any filter $G \subseteq \mathbb{P}$, let  $C_G \coloneqq \bigcup_{p \in G} C_p$. We first claim that, for any filter $G$, $C_G$ is a meet-subsemilattice of $L$ whose elements have at most $n+1$ lower covers.

\begin{claim}
Let $G\subseteq \mathbb{P}$ be a filter. Then, every element of $C_G$ has at most $n+1$ lower covers in $C_G$.
\end{claim}
\begin{proof}
Pick $x \in C_G$ and let $\vec{\alpha} = \vec{\alpha}(x)$. It suffices to prove that there exists a set $S \subseteq C_G$ of size at most $n+1$ such that $z < x$ for all $z \in S$ and, for all $y \in C_G$ with $y < x$, there is some $z \in S$ with $y \le z$.

First let 
\[
S' \coloneqq \big\{\pi_{(\vec{\alpha} \upharpoonright i)^\smallfrown (\alpha_i - 1)}(x) : i < n \text{ and } \alpha_i > 0\big\}.
\]
Now notice that, by (b), the set $C_G \cap J_{\vec{\alpha}}$ is a chain in $(L, \le)$. If there is $c \in C_G \cap J_{\vec{\alpha}}$ with $c < x$, let $c_0$ be the $\le$-greatest such $c$ and let $S = S' \cup \{c_0\}$; otherwise let $S = S'$.

We claim that the set $S$ satisfies the desired property. Observe first that $S \subseteq C_G$. Indeed, if $c_0$ is defined, then $c_0 \in C_G$ by definition. The other elements of $S$ belong to $C_G$ by (c). Moreover, all the elements of $S$ are strictly below $x$. The element $c_0$ is strictly below $x$ by definition. The other elements also cannot coincide with $x$ because, by the $\le_{lex}$-minimality of $\vec{\alpha}$, $x \not\in I_{(\vec{\alpha} \upharpoonright i)^\smallfrown (\alpha_i - 1)}$ for every  $i < n$ with $\alpha_i > 0$.

Now pick any $y \in C_G$ with $y < x$, towards showing that $y \le z$ for some $z \in S$. Let $\vec{\beta} = \vec{\alpha}(y)$. We already noted that $\vec{\beta} \le_{lex} \vec{\alpha}$. 

If $\vec{\beta} = \vec{\alpha}$, then $c_0$ is defined and $y \le c_0 \in S$ by definition of $c_0$. If $\vec{\beta} <_{lex} \vec{\alpha}$, let $i$ be such that $\vec{\beta} \upharpoonright i = \vec{\alpha} \upharpoonright i$ and $\beta_i < \alpha_i$. By (2), $I_{\vec{\beta}} \subseteq I_{\vec{\beta} \upharpoonright (i+1)}$. By (3), $I_{\vec{\beta} \upharpoonright (i+1)} \subseteq I_{(\vec{\alpha} \upharpoonright i)^\smallfrown (\alpha_i-1)}$. Thus, $y \in I_{(\vec{\alpha} \upharpoonright i)^\smallfrown (\alpha_i-1)}$. We conclude $y \le \pi_{(\vec{\alpha} \upharpoonright i)^\smallfrown (\alpha_i-1)}(x) \in S$, and we are done.
\end{proof}

\begin{claim}
Let $G\subseteq \mathbb{P}$ be a filter. Then, $C_G$ is a meet-subsemilattice of $L$.
\end{claim}
\begin{proof}
Fix $x,y \in C_G$. Let $\vec{\alpha} = \vec{\alpha}(x \wedge y)$.  By (c), both $\pi_{\vec{\alpha}} (x)$ and $\pi_{\vec{\alpha}} (y)$ belong to $C_G$. Moreover, by Lemma~\ref{lemma:meetproj}, $x \wedge y = \pi_{\vec{\alpha}} (x) \wedge \pi_{\vec{\alpha}} (y)$. In particular, $ \pi_{\vec{\alpha}}(x)$ and $ \pi_{\vec{\alpha}}(y)$ actually lie in $J_{\vec{\alpha}}$, not just in $I_{\vec{\alpha}}$.

As we already noted, the set $C_G \cap J_{\vec{\alpha}}$ is a chain. Thus, $\pi_{\vec{\alpha}}(x)$ and $\pi_{\vec{\alpha}}(y)$ must be comparable. But this means that $x \wedge y = \pi_{\vec{\alpha}}(x) \wedge \pi_{\vec{\alpha}}(y) \in \pi_{\vec{\alpha}}``\{x,y\} \subseteq C_G$. Therefore, $C_G$ is a meet-subsemilattice of $L$.
\end{proof}
  
Now we prove that if $G\subseteq \mathbb{P}$ is a $V$-generic filter, the set $C_G$ is also cofinal in $L$. Note that, by Remark~\ref{rmk:cofinalmeet}, this implies that $C_G$ is an $(n+1)$-ladder in the induced ordering.
\begin{claim}
Let $G\subseteq \mathbb{P}$ be a $V$-generic filter. Then, $C_G$ is cofinal in $L$.
\end{claim}
\begin{proof}
Fix $x \in L$ and $p \in \mathbb{P}$. It suffices by density to show that there are $q \le p$ and $y \in C_q$ with $x \le y$.  

First pick any $p' \le p$ such that, for each $(\vec{\alpha}, m) \in \mathrm{dom}(p')$, $(\vec{\alpha}, k) \in \mathrm{dom}(p')$ for every $k \le m$. 
Let
\[
y \coloneqq x \vee \bigvee \big\{p' (\vec{\alpha}, m) : (\vec{\alpha}, m) \in \mathrm{dom}(p')\big\}.
\]
Let $R = \{\vec{\alpha}(z) : z \le y\}$. Since $L$ is lower finite, $R$ is finite. For every $\vec{\alpha} \in R$, let $m_{\vec{\alpha}}$ be the least $m \in\omega$ such that $(\vec{\alpha}, m) \not\in \mathrm{dom}(p')$. Finally, let 
\[
q \coloneqq p' \cup \big\{((\vec{\alpha}, m_{\vec{\alpha}}), \pi_{\vec{\alpha}}(y)) : \vec{\alpha} \in R\big\}.
\]
Clearly, $q \le p' \le p$. We are done once we prove $y \in C_q$.  We do so by showing that $\pi_{\vec{\alpha}}(y) \in C_q$ for all $\vec{\alpha} \in R$ by induction on the lexicographic ordering of $R$---indeed, $\vec{\alpha}(y) \in R$ and $\pi_{\vec{\alpha}(y)}(y) = y$. %Note that $\vec{\alpha}(\pi_{\vec{\alpha}}(y)) = \vec{\alpha}$ for all $\vec{\alpha} \in R$.

 Fix $\vec{\alpha} \in R$ and suppose that $\pi_{\vec{\beta}} (y) \in C_q$ for all $\vec{\beta} \in R$ with $\vec{\beta} <_{lex} \vec{\alpha}$. We want to prove that $\pi_{\vec{\alpha}} (y) \in C_q$. In particular, we show that $\pi_{\vec{\alpha}} (y)$ satisfies (a)-(c) in the definition of $C_q \cap J_{\vec{\alpha}}$ for $m = m_{\vec{\alpha}}$.
 
 By definition of $q$, $q(\vec{\alpha}, m_{\vec{\alpha}}) = \pi_{\vec{\alpha}}(y)$. In particular, $\pi_{\vec{\alpha}}(y)$ satisfies (a). Moreover, since  $y \ge q(\vec{\alpha}, k) = p'(\vec{\alpha}, k) $ for all $k  < m_{\vec{\alpha}}$ by definition of $y$, we conclude that $\pi_{\vec{\alpha}}(y) \ge q(\vec{\alpha}, k)$ for all $k  < m_{\vec{\alpha}}$. Thus, $\pi_{\vec{\alpha}}(y)$ also satisfies (b). 
 
Finally, towards showing that $\pi_{\vec{\alpha}}(y)$ satisfies also (c), pick $\vec{\gamma} \in T$ with $\vec{\gamma} <^*_{lex} \vec{\alpha}$. We need to prove $\pi_{\vec{\gamma}} \circ \pi_{\vec{\alpha}}(y) \in C_q$. Denote $\pi_{\vec{\gamma}} \circ \pi_{\vec{\alpha}}(y) $ by $z$. Then
\begin{equation}\label{eq:finalcohen}
\begin{split}
z &= \pi_{\vec{\alpha}(z)} \circ \pi_{\vec{\gamma}} \circ \pi_{\vec{\alpha}} (y)\\
&= \pi_{\vec{\alpha}(z)} \circ \pi_{\vec{\alpha}} (y)\\
&= \pi_{\vec{\alpha}(z)}(y).
\end{split}
\end{equation}
The first equality follows trivially, as $z = \pi_{\vec{\alpha}(z)} (z)$. To argue the second equality, note that $I_{\vec{\alpha}(z)} \subseteq I_{\vec{\gamma}}$. Indeed, there are two cases: either $\vec{\gamma} \subseteq \vec{\alpha}(z)$, and in this case $I_{\vec{\alpha}(z)} \subseteq I_{\vec{\gamma}}$ by (2); or $\vec{\alpha}(z) <_{lex} \vec{\gamma}$ and then by (6) we conclude $I_{\vec{\alpha}(z)} \subseteq I_{\vec{\gamma}}$. So in either case $I_{\vec{\alpha}(z)} \subseteq I_{\vec{\gamma}}$. Then, $\pi_{\vec{\alpha}(z)} \circ \pi_{\vec{\gamma}} = \pi_{\vec{\alpha}(z)}$ follows from Lemma~\ref{lemma:concproj} since $I_{\vec{\alpha}(z)} \subseteq I_{\vec{\gamma}}$. Finally, the third equality also follows from Lemma~\ref{lemma:concproj}. Indeed, $\vec{\alpha}(z)  <_{lex} \vec{\alpha}$, and as such we have $I_{\vec{\alpha}(z)} \subseteq I_{\vec{\alpha}}$ by (6).  

Thus, from \eqref{eq:finalcohen} we conclude that $z = \pi_{\vec{\alpha}(z)} (y)$. But $\vec{\alpha}(z) \in R$ and $\vec{\alpha}(z) <_{lex} \vec{\alpha}$, and thus by induction hypothesis $\pi_{\vec{\alpha}(z)} (y) \in C_q$. We conclude $z \in C_q$ and, overall, that $\pi_{\vec{\alpha}}(y)$ satisfies also (c).
\end{proof}
\end{proof}

\end{document}
